\section{Deep Generative Models for Medical Imaging}
\label{sec:medical_imaging_context}
Deep generative models have transformed how imaging problems are formulated and addressed in biomedical research. 
Unlike traditional deterministic frameworks that rely on handcrafted priors or model-based regularization, 
modern generative paradigms employ learned probabilistic representations to synthesize, reconstruct, and infer 
data under uncertainty \citep{diffusion_survey_2025}. 
These methods differ primarily in how they represent and learn the data distribution \(P_{\text{data}}(x)\), 
ranging from adversarial and likelihood-based formulations to diffusion and energy-based perspectives. 
Figure~\ref{fig:deep_generative_models_overview} summarises the principal categories of deep generative models 
and their core mechanisms, illustrating how each family balances interpretability, likelihood tractability, and sample fidelity 
across diverse applications in medical image analysis.

\begin{figure}[H]
    \centering
    \includegraphics[width=0.85\textwidth]{chapters/background/media/medical_imaging/medical_examples_of_dgm.png}
    \caption{
        \textbf{Overview of Deep Generative Models \citep{kazerouni2022diffusion}.}
        Major families of generative models include (a) Generative Adversarial Networks (GANs), which train a generator and discriminator adversarially; 
        (b) Energy-Based Models (EBMs), which learn an unnormalized energy landscape over data; 
        (c) Variational Autoencoders (VAEs), which encode data into a latent distribution and decode samples from it; 
        (d) Normalizing Flows (NFs), which learn invertible mappings for exact likelihood computation; 
        and (e) Diffusion Models, which progressively corrupt and denoise data through a forward–reverse stochastic process.
    }
    \label{fig:deep_generative_models_overview}
\end{figure}



\subsection{Generative Formulations in Medical Imaging}

Medical imaging problems can be broadly categorized by the objective that generative modeling aims to address. 
These include reconstruction-oriented tasks, inverse imaging problems, image translation, denoising and restoration, 
data augmentation, and recently, compositional or prior-based inference frameworks. 

\subsubsection{Inverse Imaging and Reconstruction Problems}
Inverse problems in biomedical imaging \citep{accelerated_mri_motion_correction_2023, score_based_diffusion_mri_2022, cao2022spiritdiffusion}, often applied in MRI reconstruction, sparse-view CT, and low-dose denoising, arise when the observed measurements \( y \) are an indirect or incomplete observation of the true physical image \( x \). 
Formally, such problems are described by the forward model:
\[
y = A(x) + \epsilon,
\]
where \( A(\cdot) \) denotes the measurement or imaging operator (e.g., Fourier transform in MRI, Radon transform in CT, 
or projection operator in X-ray imaging), and \( \epsilon \) denotes measurement noise.

\paragraph{Generative Perspective.}
Generative models are used as data-driven priors, replacing handcrafted regularizers. 
A Bayesian reconstruction is then written as:
\[
p(x|y) \propto p(y|x)\,p_\theta(x),
\]
where \( p_\theta(x) \) is a learned generative prior such as a diffusion or score-based model. 
This allows reconstructions that are both \emph{data-consistent} and \emph{anatomically plausible}. 

\subsubsection{Image-to-Image Translation and Synthesis}
In Image-to-Image translation tasks, generative models such as GANs, VAEs, and diffusion models are used to translate between domains \citep{Lyu2022ConversionBC, Meng2022ANU}:
\[
x_\text{target} = G_\theta(x_\text{source}),
\]
where \( G_\theta \) is a learned mapping from source modality \( x_\text{source} \) to target modality \( x_\text{target} \).
Common applications include cross-modality synthesis, for example, synthesizing CT images from MR, or generating contrast-enhanced from non-contrast scans. 
These approaches rely on learned joint distributions \( p_\theta(x_\text{target}, x_\text{source}) \) and enable 
cross-modality synthesis for diagnostic or data augmentation purposes.

\subsubsection{Data Augmentation and Distribution Learning}
A major use case of generative models is to learn the data distribution \( p_\theta(x) \approx p_\text{data}(x) \) 
and sample from it to generate realistic synthetic examples:
\[
x \sim p_\theta(x).
\]
Such models are used to balance class distributions, enhance rare-pathology representation, or augment limited datasets \citep{generative_ai_medical_imaging_2025}. 
Diffusion models have shown remarkable fidelity in generating realistic medical images, including chest radiographs, 
brain MRIs, and retinal scans.

\subsection{Problem Categories for Chest Radiography}

Within this taxonomy, chest radiographs occupy a distinctive position. 
They are typically \emph{fully reconstructed 2D projections} of thoracic anatomy, 
so classical reconstruction-based inverse problems (as seen in MRI or CT) are generally not applicable to standard X-rays.
Instead, generative modeling for chest radiography primarily addresses \citep{generative_ai_medical_imaging_2025}:
\begin{itemize}
    \item \textbf{Data augmentation:} generating realistic X-rays for rare diseases such as tuberculosis (TB) or silicosis;
    \item \textbf{Denoising and artifact removal:} treating noise and motion artifacts as mild inverse problems;
    \item \textbf{Representation learning:} capturing latent manifolds of disease progression and pathology;
\end{itemize}

To the best of our knowledge, there has not yet been a case study where researchers leverage compositionality or pretrained diffusion models \emph{at inference time} to solve a problem pertinent to medical imaging. 
The work presented in this research explores precisely this space: it falls within the realm of \emph{compositional generative modeling}, where pretrained diffusion models trained on distinct chest X-ray subsets (e.g., Normal and TB) are combined at inference to produce hybrid or interpolated samples representing compositional pathology manifolds. 
This formulation can be regarded as a generalization of mixture-of-experts and guided diffusion sampling.
